GENIUS and other generative multimodal retrievers decode a discrete semantic ID -- from a residual quantizer (RQ) -- instead of dense nearest-neighbor search, making retrieval cost independent of corpus size. We reproduce GENIUS end-to-end from its officially released Stage-1 checkpoint, then study whether regularizing the RQ codebook toward balanced usage helps or hurts retrieval, via a controlled, multi-seed experiment across MSCOCO, VisualNews (broad-domain), and FashionIQ (fine-grained). On MSCOCO, strong regularization significantly improves mean text-to-image Recall@1 ($23.33\%$ vs.\ $20.78\%$, five decoder seeds, $+12.3\%$ relative, Welch's $t$-test $p<10^{-6}$), but reliably collapses image-to-text Recall from $10.30\%$ to near zero on the identical checkpoints. Since both effects could be artifacts of the single RQ tokenizer each condition was built from, we independently retrained the tokenizer twice more: across three fully independent tokenizers, every regularized one still beats every vanilla one on text-to-image ($24.32\%$ vs.\ $21.16\%$ mean, non-overlapping), and the collapse holds on all three. A dense, decoder-free probe over the frozen tokenizer's embeddings shows the same collapse before any decoder training, partially explained by text embeddings' greater anisotropy under regularization. On VisualNews, the same probe shows both effects are tokenizer-level too, but text-to-image reverses: strong regularization more than halves Recall instead of raising it, while image-to-text collapses exactly as on MSCOCO -- the collapse is architectural and generalizes, the gain is real but dataset-specific. FashionIQ is hurt by any regularization. The natural fix -- decoupling regularization per target modality, including regularizing only the text side -- fails regardless of which side carries the weight: image-to-text stays collapsed even at zero text-side weight, and text-to-image gets worse than either baseline, because image and text rows share one encoder and codebook. No single balance-regularization strategy transfers across retrieval regimes, and even the one clear positive result does not generalize beyond the dataset where we found it.
